trallala

Emmy leefde in de kerk
bedolven onder 't werk.

Vier keer per dag
exact op de klokslag

kwam een monnik
na enig gefrunnik

haar losmaken
zonder aanraken.

Gelijk een noodgeval
in de zwijnenstal

mocht zij ontlasten.
Minuten lang plasten

de overvolle blazen
meerdere bierglazen

vol over de smurrie
van pis en durrie.

Een emmer koud water
gutste ietwat later

heel aangenaam
over haar lichaam.

Weg was de kleefpasta
van kwijl en sperma.

Verder kwam de jachtbuit
er nauwelijks meer uit.

Maandelijks niettemin,
als de zon onderging,

werd ze gesnoerd
en abrupt ontvoerd,

in boeien geslagen
op een hooiwagen,

waarop een kooi stond.
Die bracht haar terstond

naar een geheime locatie;
de bisschopsresidentie.

Daar zat de Bisschop
met zijn grijze kop

gedoken in de boeken
de waarheid te zoeken.

Negerend haar intrede,
viel ze hem in de rede.

Een vinnige blik
bezorgde schrik.

Ze viel met bibberatie
meteen in prostratie.

Smekend om vergeving.
Die ze uiteindelijk ontving.

Spijts zijn hoge positie
was zijn cohabitatie

niet erg excentriek
doch beneden alle kritiek.

Hij had de eterniteit
met de delicieuze meid.

Geen jachtige wachtrij
De standenmaatschappij

gaf hem dit voordeel.
Evenzeer confessioneel

waren de colloquia
die volgden erna.

Hij zette uiteen dat
de cultuur van haar stad

tamelijk excentrisch was.
De kerk had connivance

met de praktijken.
Zo zou ras blijken.

Een grote bijdrage
aan de pauselijke gage

gaf hen de autonomie
voor een eigen religie.

De handel en wandel
leverde lucratieve handel.

De vele immigranten
werden declaranten

mits hun stemrecht
werd toegezegd

in het stadsdebat
over de geliefde stad.

Met het burgerrecht
konden schurken terecht

en zich vrijkopen
om straf te ontlopen.

Een vrouw geschaakt
nam men naakt

mee de stad in.
Mogelijke tegenzin

werd dan miskent.
Als bezit toegekend.

Men had dan het recht
te verbinden in de echt

Werd dit nagelaten
dan moest men de baten

delen met de parochie
tijdens de communie.

De talrijke vrijgezellen
mochten het stellen

met de buurvrouw
als die dat ook wou.

Net als de echtgenoot
die vrijelijk genoot

van het beminnen
van de lustslavinnen

die de parochie beheerde
en dagelijks exhibeerde.

Iedere stadjer was gedrild
“Heb lief en doe wat je wilt”.

Zo gebruikten de parvenu’s
Augustinus’ filosofie als excuus.

Deze corrupte Bisschop
van een stad verderop

misbruikte deze kennis
voor heiligschennis.

Hij eiste elke maand
dat er, tegenoverstaand

dat er niet werd gestraft,
een snol werd verschaft

als solidariteitstoeslag
voor minstens een dag.

Het celibaat was zwaar.
Dus was het niet raar

dat de bisschop deelde
met een minderbedeelde.

Zo verkreeg hij ruggesteun
van een andere miskleun.

Het aldoende gade slaan
mocht hem niet ontgaan

Als Emmy moest afzuigen
kon hij ervan getuigen

dat hij haar poesje vulde.
Waarna hij onthulde

hetgeen in haar geslopen
eruit wist te lopen.

Emmy welgeleerd
lag snel omgekeerd.
